\section{Introduction}
\label{sec:intro}

% Framing: the field's evidentiary problem, not just "we propose a model".
\paragraph{Motivation.}
Hybrid quantum--classical models are increasingly proposed for financial
prediction, and the reported advantages are typically large: the reference
study this work audits reports a Sharpe ratio \SI{72}{\percent} above its best
classical baseline \citep{chen2025quantumb}. Yet the applied literature
shares three vulnerabilities --- short evaluation samples, universes assembled
with hindsight, and weak baselines --- and controlled designs that fix
everything but the quantum component find no advantage on the equity
cross-section \citep{shen2026quantum,kobayashi2023cross}. Benchmarks outside
finance reach the same conclusion: prototypical classical baselines beat
prototypical quantum models, entanglement can be removed at no cost
\citep{bowles2024better}, and where a quantum layer appears to help, a
parameter-matched classical control usually helps as much
\citep{gillani2026quantum}. At the same time, the finance literature has
developed tools --- the Deflated Sharpe Ratio, the probability of backtest
overfitting, multiplicity-robust inference, decision-time leakage audits ---
that the quantum finance literature has not yet adopted; no study in our
literature search reports any of them for a quantum model.

\paragraph{This is a self-audit.}
The reference study whose result this paper re-examines is our own
\citep{chen2025quantumb}. We state that at the outset because it determines how
the rest should be read. The protocol below was not assembled to test someone
else's claim; it was assembled to test a claim we had already published, and it
returned the opposite answer. Nothing in the analysis depends on that fact ---
the conditions, the windows and the tests are specified in advance of the
result and the data and code are released --- but a reader is entitled to know
that the paper's critical findings are directed at its own authors' prior work
rather than at a third party's. A literature search of \num{2107} deduplicated
records for this paper found a single study that drew critical conclusions from
its own empirical results \citep{bowles2024better}; we would rather add to that
count than to the other one.

\paragraph{The identification problem.}
A hybrid model typically routes both a quantum score and the classical encoder
representation into a shared head. When such a model outperforms a baseline, the
result is consistent with at least four explanations: the circuit learned a
useful nonlinearity; the circuit acts as a fixed random feature; the classical
encoder did all the work; or the difference is sampling noise. Distinguishing
them requires control conditions, not a larger table of baselines.

\paragraph{Contributions.}
\begin{enumerate}
  \item A controlled ablation that isolates the quantum contribution into four
    conditions --- trained circuit, frozen circuit, parameter-matched classical
    surrogate, and no circuit --- trained through an identical code path.
  \item An evaluation protocol for cross-sectional ranking that reports
    annualised metrics on \num{500} walk-forward out-of-sample days with
    bootstrap confidence intervals, Deflated Sharpe Ratios, backtest-overfitting
    probabilities and multiple-comparison control.
  \item A batched formulation of the quantum convolutional layer that reproduces
    the per-sample implementation exactly in double precision while running
    \num{33} to \num{72} times faster depending on batch size, which is what
    makes full-sample walk-forward evaluation tractable.
  \item A re-evaluation of our own previously published result
    \citep{chen2025quantumb} on \num{241} out-of-sample days rather than the
    \num{22} it reported, which decomposes the advantage it claimed into four
    identifiable sources, one of which is a split adjustment applied in the
    wrong temporal direction that leaves more artefactual returns than
    performing no adjustment at all.
\end{enumerate}
